% GENERATED FILE. Edit canonical lessons, then run npm run generate:books.
% canonical-source-hash: fnv1a64:72cc6ba55d5b04d7
% canonical-lessons: HI-C511-cak, HI-C511-sanna, HI-C511-ava, HI-C511-saca, HI-C511-dhaalna

\chapter{Clay and Fire}
\label{ch:hi-clay-and-fire}
\hlchaptermodality{511}

\begin{chapteropening}
\textbf{By the end of this chapter:} \emph{I can talk about how a potter works.}

Complete the last lesson of chapter 511: I can talk about how a potter works.
\end{chapteropening}

\section[\texorpdfstring{\hi{चाक} (cāk)}{चाक (cāk)}]{\hi{चाक} (cāk) — a potter's wheel}

\label{lesson:HI-C511-cak}



\begin{quote}
Before the new one: say the Hindi for to knock in, then the Hindi for to bore.
\end{quote}

\subsection*{You'll want to know: \hi{चाक}}
\textbf{\hi{चाक}} — \emph{cāk} — \textquotedblleft{}a potter's wheel\textquotedblright{}.

\begin{quote}
\hi{कुम्हार} \hi{चाक} \hi{पर} \hi{मिट्टी} \hi{के} \hi{दीये} \hi{बनाता} \hi{है।}
\end{quote}

\begin{quote}
\emph{Kumhār cāk par miṭṭī ke dīye banātā hai.}
\end{quote}

\begin{quote}
The potter makes clay lamps on his wheel.
\end{quote}

Masculine.

\begin{cousinweb}
\textbf{\hi{चाक}} (\emph{cāk}) is a tadbhava, inherited from Sanskrit \emph{cakra}, \textquotedblleft{}a wheel\textquotedblright{} (CDIAL 4538). CDIAL gives the Hindi word as \textquotedblleft{}any kind of wheel\textquotedblright{}; at the potter's it is his wheel.
\end{cousinweb}

\begin{grammarlens}[title={What you've built}]
One more word for describing what you lived and felt.
\end{grammarlens}

\subsection*{Your turn}
\begin{itemize}
\raggedright
  \item \emph{Say it:} \emph{cāk}
  \item \emph{Say it:} \emph{cāk}, once more
  \item \emph{Say it:} \emph{Kumhār cāk par miṭṭī ke dīye banātā hai.} — The potter makes clay lamps on his wheel.
  \item \emph{Recall:} say the Hindi for to knock in, then the Hindi for to bore, then say \emph{cāk} again
\end{itemize}

\begin{tcolorbox}[breakable,colback=teal!4,colframe=teal!35!black,title={Before you move on}]
What does \hi{चाक} mean? (\textquotedblleft{}A potter's wheel\textquotedblright{}.) Say it once more.
\end{tcolorbox}
\section[\texorpdfstring{\hi{सानना} (sānnā)}{सानना (sānnā)}]{\hi{सानना} (sānnā) — to mix (with water); to knead (clay)}

\label{lesson:HI-C511-sanna}



\begin{quote}
Before the new one: say the Hindi for to bore, then the Hindi for a potter's wheel.
\end{quote}

\subsection*{You'll want to know: \hi{सानना}}
\textbf{\hi{सानना}} — \emph{sānnā} — \textquotedblleft{}to mix (with water); to knead (clay)\textquotedblright{}.

\begin{quote}
\hi{कुम्हार} \hi{मिट्टी} \hi{को} \hi{पानी} \hi{में} \hi{सानता} \hi{है।}
\end{quote}

\begin{quote}
\emph{Kumhār miṭṭī ko pānī meṁ sāntā hai.}
\end{quote}

\begin{quote}
The potter mixes the clay with water.
\end{quote}

Past, with \hi{ने} (he or she …): \hi{उसने} \hi{मिट्टी} \hi{सानी}, \emph{usne miṭṭī sānī}.

\begin{cousinweb}
\textbf{\hi{सानना}} (\emph{sānnā}) is a tadbhava, inherited from Sanskrit \emph{saṁdadhāti}, \textquotedblleft{}places on, combines\textquotedblright{} (CDIAL 12898). CDIAL gives the Hindi verb as \textquotedblleft{}to mix, knead\textquotedblright{}.
\end{cousinweb}

\begin{grammarlens}[title={What you've built}]
One more word for describing what you lived and felt.
\end{grammarlens}

\subsection*{Your turn}
\begin{itemize}
\raggedright
  \item \emph{Say it:} \emph{sānnā}
  \item \emph{Say it:} \emph{sānnā}, once more
  \item \emph{Say it:} \emph{Kumhār miṭṭī ko pānī meṁ sāntā hai.} — The potter mixes the clay with water.
  \item \emph{Recall:} say the Hindi for to bore, then the Hindi for a potter's wheel, then say \emph{sānnā} again
\end{itemize}

\begin{tcolorbox}[breakable,colback=teal!4,colframe=teal!35!black,title={Before you move on}]
What does \hi{सानना} mean? (\textquotedblleft{}To mix (with water); to knead (clay)\textquotedblright{}.) Say it once more.
\end{tcolorbox}
\section[\texorpdfstring{\hi{आँवा} (ā̃vā)}{आँवा (ā̃vā)}]{\hi{आँवा} (ā̃vā) — a potter's kiln}

\label{lesson:HI-C511-ava}



\begin{quote}
Before the new one: say the Hindi for a potter's wheel, then the Hindi for to mix (with water).
\end{quote}

\subsection*{You'll want to know: \hi{आँवा}}
\textbf{\hi{आँवा}} — \emph{ā̃vā} — \textquotedblleft{}a potter's kiln\textquotedblright{}.

\begin{quote}
\hi{कुम्हार} \hi{ने} \hi{कच्चे} \hi{बर्तन} \hi{आँवे} \hi{में} \hi{रखे।}
\end{quote}

\begin{quote}
\emph{Kumhār ne kacce bartan ā̃ve meṁ rakhe.}
\end{quote}

\begin{quote}
The potter put the unfired pots in the kiln.
\end{quote}

Masculine.

\begin{cousinweb}
\textbf{\hi{आँवा}} (\emph{ā̃vā}) is a tadbhava, inherited from Sanskrit \emph{āpāka}, \textquotedblleft{}a potter's kiln\textquotedblright{} (CDIAL 1202).
\end{cousinweb}

\begin{grammarlens}[title={What you've built}]
One more word for describing what you lived and felt.
\end{grammarlens}

\subsection*{Your turn}
\begin{itemize}
\raggedright
  \item \emph{Say it:} \emph{ā̃vā}
  \item \emph{Say it:} \emph{ā̃vā}, once more
  \item \emph{Say it:} \emph{Kumhār ne kacce bartan ā̃ve meṁ rakhe.} — The potter put the unfired pots in the kiln.
  \item \emph{Recall:} say the Hindi for a potter's wheel, then the Hindi for to mix (with water), then say \emph{ā̃vā} again
\end{itemize}

\begin{tcolorbox}[breakable,colback=teal!4,colframe=teal!35!black,title={Before you move on}]
What does \hi{आँवा} mean? (\textquotedblleft{}A potter's kiln\textquotedblright{}.) Say it once more.
\end{tcolorbox}
\section[\texorpdfstring{\hi{साँचा} (sā̃cā)}{साँचा (sā̃cā)}]{\hi{साँचा} (sā̃cā) — a mould}

\label{lesson:HI-C511-saca}



\begin{quote}
Before the new one: say the Hindi for to mix (with water), then the Hindi for a potter's kiln.
\end{quote}

\subsection*{You'll want to know: \hi{साँचा}}
\textbf{\hi{साँचा}} — \emph{sā̃cā} — \textquotedblleft{}a mould\textquotedblright{}.

\begin{quote}
\hi{वह} \hi{मिट्टी} \hi{को} \hi{साँचे} \hi{में} \hi{दबाकर} \hi{खिलौने} \hi{बनाती} \hi{है।}
\end{quote}

\begin{quote}
\emph{Vah miṭṭī ko sā̃ce meṁ dabākar khilaune banātī hai.}
\end{quote}

\begin{quote}
She makes toys by pressing clay into a mould.
\end{quote}

Masculine.

\begin{cousinweb}
\textbf{\hi{साँचा}} (\emph{sā̃cā}) is a tadbhava, inherited from Sanskrit \emph{sañcaka}, \textquotedblleft{}a mould\textquotedblright{} (CDIAL 13096).
\end{cousinweb}

\begin{grammarlens}[title={What you've built}]
One more word for describing what you lived and felt.
\end{grammarlens}

\subsection*{Your turn}
\begin{itemize}
\raggedright
  \item \emph{Say it:} \emph{sā̃cā}
  \item \emph{Say it:} \emph{sā̃cā}, once more
  \item \emph{Say it:} \emph{Vah miṭṭī ko sā̃ce meṁ dabākar khilaune banātī hai.} — She makes toys by pressing clay into a mould.
  \item \emph{Recall:} say the Hindi for to mix (with water), then the Hindi for a potter's kiln, then say \emph{sā̃cā} again
\end{itemize}

\begin{tcolorbox}[breakable,colback=teal!4,colframe=teal!35!black,title={Before you move on}]
What does \hi{साँचा} mean? (\textquotedblleft{}A mould\textquotedblright{}.) Say it once more.
\end{tcolorbox}
\section[\texorpdfstring{\hi{ढालना} (ḍhālnā)}{ढालना (ḍhālnā)}]{\hi{ढालना} (ḍhālnā) — to pour into a mould; to cast}

\label{lesson:HI-C511-dhaalna}



\begin{quote}
Before the new one: say the Hindi for a potter's kiln, then the Hindi for a mould.
\end{quote}

\subsection*{You'll want to know: \hi{ढालना}}
\textbf{\hi{ढालना}} — \emph{ḍhālnā} — \textquotedblleft{}to pour into a mould; to cast\textquotedblright{}.

\begin{quote}
\hi{लोहार} \hi{ने} \hi{गरम} \hi{लोहा} \hi{साँचे} \hi{में} \hi{ढाला।}
\end{quote}

\begin{quote}
\emph{Lohār ne garam lohā sā̃ce meṁ ḍhālā.}
\end{quote}

\begin{quote}
The blacksmith poured hot iron into the mould.
\end{quote}

Past, with \hi{ने} (he or she …): \hi{उसने} \hi{एक} \hi{घंटी} \hi{ढाली}, \emph{usne ek ghaṇṭī ḍhālī}.

\begin{cousinweb}
\textbf{\hi{ढालना}} (\emph{ḍhālnā}) is a tadbhava, inherited from \emph{*ḍhalati}, \textquotedblleft{}bends over, falls\textquotedblright{}, a form CDIAL reconstructs from its descendants (CDIAL 5581). CDIAL lists Hindi \emph{ḍhālnā}, \textquotedblleft{}to pour\textquotedblright{}, under its causative; casting is pouring into a mould.
\end{cousinweb}

\begin{grammarlens}[title={What you've built}]
One more word for describing what you lived and felt.
\end{grammarlens}

\subsection*{Your turn}
\begin{itemize}
\raggedright
  \item \emph{Say it:} \emph{ḍhālnā}
  \item \emph{Say it:} \emph{ḍhālnā}, once more
  \item \emph{Say it:} \emph{Lohār ne garam lohā sā̃ce meṁ ḍhālā.} — The blacksmith poured hot iron into the mould.
  \item \emph{Recall:} say the Hindi for a potter's kiln, then the Hindi for a mould, then say \emph{ḍhālnā} again
\end{itemize}

\begin{tcolorbox}[breakable,colback=teal!4,colframe=teal!35!black,title={Before you move on}]
What does \hi{ढालना} mean? (\textquotedblleft{}To pour into a mould; to cast\textquotedblright{}.) Say it once more.
\end{tcolorbox}
